% Options for packages loaded elsewhere
\PassOptionsToPackage{unicode}{hyperref}
\PassOptionsToPackage{hyphens}{url}
%
\documentclass[
]{article}
\title{Christian Hip-Hop and Moral Priviledge\thanks{Thanks be to the
ASA. \textbf{Current version}: November 30, 2024; \textbf{Corresponding
author}:
\href{mailto:montaque.reynolds@slu.edu}{\nolinkurl{montaque.reynolds@slu.edu}}.}}
\author{true}
\date{}

\usepackage{amsmath,amssymb}
\usepackage{lmodern}
\usepackage{iftex}
\ifPDFTeX
  \usepackage[T1]{fontenc}
  \usepackage[utf8]{inputenc}
  \usepackage{textcomp} % provide euro and other symbols
\else % if luatex or xetex
  \usepackage{unicode-math}
  \defaultfontfeatures{Scale=MatchLowercase}
  \defaultfontfeatures[\rmfamily]{Ligatures=TeX,Scale=1}
\fi
% Use upquote if available, for straight quotes in verbatim environments
\IfFileExists{upquote.sty}{\usepackage{upquote}}{}
\IfFileExists{microtype.sty}{% use microtype if available
  \usepackage[]{microtype}
  \UseMicrotypeSet[protrusion]{basicmath} % disable protrusion for tt fonts
}{}
\makeatletter
\@ifundefined{KOMAClassName}{% if non-KOMA class
  \IfFileExists{parskip.sty}{%
    \usepackage{parskip}
  }{% else
    \setlength{\parindent}{0pt}
    \setlength{\parskip}{6pt plus 2pt minus 1pt}}
}{% if KOMA class
  \KOMAoptions{parskip=half}}
\makeatother
\usepackage{xcolor}
\IfFileExists{xurl.sty}{\usepackage{xurl}}{} % add URL line breaks if available
\IfFileExists{bookmark.sty}{\usepackage{bookmark}}{\usepackage{hyperref}}
\hypersetup{
  pdftitle={Christian Hip-Hop and Moral Priviledge},
  pdfkeywords={Care Ethics, Aesthetics},
  hidelinks,
  pdfcreator={LaTeX via pandoc}}
\urlstyle{same} % disable monospaced font for URLs
\usepackage[margin=1in]{geometry}
\usepackage{graphicx}
\makeatletter
\def\maxwidth{\ifdim\Gin@nat@width>\linewidth\linewidth\else\Gin@nat@width\fi}
\def\maxheight{\ifdim\Gin@nat@height>\textheight\textheight\else\Gin@nat@height\fi}
\makeatother
% Scale images if necessary, so that they will not overflow the page
% margins by default, and it is still possible to overwrite the defaults
% using explicit options in \includegraphics[width, height, ...]{}
\setkeys{Gin}{width=\maxwidth,height=\maxheight,keepaspectratio}
% Set default figure placement to htbp
\makeatletter
\def\fps@figure{htbp}
\makeatother
\setlength{\emergencystretch}{3em} % prevent overfull lines
\providecommand{\tightlist}{%
  \setlength{\itemsep}{0pt}\setlength{\parskip}{0pt}}
\setcounter{secnumdepth}{-\maxdimen} % remove section numbering
\newlength{\cslhangindent}
\setlength{\cslhangindent}{1.5em}
\newlength{\csllabelwidth}
\setlength{\csllabelwidth}{3em}
\newlength{\cslentryspacingunit} % times entry-spacing
\setlength{\cslentryspacingunit}{\parskip}
\newenvironment{CSLReferences}[2] % #1 hanging-ident, #2 entry spacing
 {% don't indent paragraphs
  \setlength{\parindent}{0pt}
  % turn on hanging indent if param 1 is 1
  \ifodd #1
  \let\oldpar\par
  \def\par{\hangindent=\cslhangindent\oldpar}
  \fi
  % set entry spacing
  \setlength{\parskip}{#2\cslentryspacingunit}
 }%
 {}
\usepackage{calc}
\newcommand{\CSLBlock}[1]{#1\hfill\break}
\newcommand{\CSLLeftMargin}[1]{\parbox[t]{\csllabelwidth}{#1}}
\newcommand{\CSLRightInline}[1]{\parbox[t]{\linewidth - \csllabelwidth}{#1}\break}
\newcommand{\CSLIndent}[1]{\hspace{\cslhangindent}#1}
\usepackage{longtable}
\usepackage{hyperref}
\usepackage{setspace}
\usepackage{amsthm}
\newtheorem{theorem}{hyp}
\newtheorem{lemma}{Lemma}
\newtheorem{theoremc}{Continuity}
\newtheorem{theoremdc}{Discontinuity}
\newtheorem{theoremcr}{Contextual Reasons}
\newtheorem{theoremfr}{Functioned Reasons}
\newtheorem{theoremfe}{Framing Effects}
\usepackage{booktabs}
\usepackage{longtable}
\usepackage{array}
\usepackage{multirow}
\usepackage{wrapfig}
\usepackage{float}
\usepackage{colortbl}
\usepackage{pdflscape}
\usepackage{tabu}
\usepackage{threeparttable}
\usepackage{threeparttablex}
\usepackage[normalem]{ulem}
\usepackage{makecell}
\usepackage{xcolor}
\ifLuaTeX
  \usepackage{selnolig}  % disable illegal ligatures
\fi

\begin{document}
\maketitle
\begin{abstract}
Narrative thinking is substantial in our lives as it provides ways of
thinking about how our lives may have gone differently, our plans for
the future, and how things might turn out (non-actual episodes). Looking
at Christian Hip-Hop, I argue that it defines an alternative narrative.
I argue that it is a moral technology that has ethical implications in a
cultural context.
\end{abstract}

\hypertarget{introduction}{%
\section{Introduction}\label{introduction}}

It is often taken for granted that themes in hip-hop are either
synonymous with issues of social justice, or misogyny. Further, it is an
open and controversial question whether or not the values of the music
communities hip-hop represents are represented in the common portrayal
of commercial hip-hop or not (Shusterman 2007, 1995; \textbf{ruth01?}).
The former is often referenced as a response to historical and
contemporary factors, such as mass incarceration, over / under policing,
and the life expectancies of communities often associated with hip-hop.
In the first case, hip-hop is important because of its social
commentary. While it may be argued that the latter themes undermine the
former, critics have often overlooked what may appear to be the most
obvious implications of these conflicting themes. Instead, critics argue
that the latter themes are either clever social commentaries in
themselves, or that profit is the primary motivation for the latter
themes.

The meaning and associations of these themes are often subjects for
intense debates from which which there is not much hope for agreement.
Here, I will consider a class of hip-hop artists that do not receive
enough attention. I will analyze this class of hip-hop from within the
context of a narrative response to reductionism. Here, Christian hip-hop
is a narrative which introduces a practical identity (\textbf{kors89?})
that offers reasons for action (See \textbf{gold08?}; \textbf{vell05?};
\textbf{alth22?}).

\hypertarget{literature-review}{%
\section{Literature Review}\label{literature-review}}

According to narrativists, narratives are an important function for
explaining an agent's motivation to perform some action. One implication
of this view then, is that agency requires one to be able to determine
her own narratives.

\begin{quote}
In understanding ourselves and others, we do so through narratives that
explain the meaning of actions by specifying what motives could make
sense of both the action and the various psychological, physiological,
and social events that preceded it. (\textbf{alth22?})
\end{quote}

In his paper \emph{Narrative Thinking, Emotion, and Planning}, Goldie
argues that narrative is substantial in our lives as it provides ways of
thinking about how our lives may have gone differently, our plans for
the future, and how things might turn out (non-actual episodes).

Christine Korsgaard offers a particular reason why this is an important
point for practical deliberation. Narratives give us ways to think about
ourselves. Who we take ourselves to be, is a fundamental factor from the
standpoint of practical reason.

\hypertarget{narrative-thinking-planning-and-agency}{%
\subsection{Narrative Thinking, Planning, and
Agency}\label{narrative-thinking-planning-and-agency}}

When we make plans for the future, our actions are often constrained. We
are constrained by facts about ourselves, for instance, we may not be
tall enough to go on a given amusement park ride. Or have the necessary
professional qualifications for a given job on the coast. Sometimes we
are constrained by contingent facts about ourselves. These can include
our personalities, our character traits, our emotional health and
professional dispositions such as being organized or a ``careful''
thinker. Each of these constraints open up particular opportunities to
us while blocking other ones. We can be motivated towards one course of
action while being discouraged from another. Some times these
constraints are external. They may be past events, for instance if we
were permanently injured in an automobile accident; they could be
current events, the economy and inflation making a vacation this year
impossible. They can be others, as when someone rear-ends us and we
become unable to go on a planned ski trip due to our injuries. We can
also be constrained by the plans that others may have for us, for
instance considering that our spouse may have always wanted to live near
her family and therefore we decline a job offer that would make this
impossible. Sometimes they are policies that others live by. Maybe our
parents would never approve of us marrying someone of a different faith.

Goldie argues that past experiences can motivate counterfactual ways of
thinking and regret. We may have shown up late to a job interview
undermining any chance we had of working for that company. Although
there are times when the constraints and experiences and the regrets
they lead to are real, occurring and constraining our choices in
real-life, Goldie argues that there are also times wherein they take
place in make-believe and other fictional events, or counterfactual and
non-local events.

Non-actual events are ones that take place in the past or in the future.
It can be a part of a fiction, or some event that occurred in
make-believe. Imagine a friend playing a prank on you. He tells you that
there is a speed-trap set up by the local law enforcement on the way you
normally take home from work. You would rather not have to worry about
receiving a potential citation. There is one other route of travel
available to you although this one takes on average twenty-five minutes
longer. A narrative is a representation of a sequence of events and your
friend has recounted a narrative to you. We become motivated to take the
alternative route in order to avoid the speed trap.

More specifically, your friend's narrative is the mental representation
of a sequence of events, whether in thought only, or in real life. The
narrative provides coherence of the causal and other connections between
events into a narrative episode or episodes. What we can glean from
these episodes are coherence, meaningfulness, and emotional
importance/import. Emotional importance of the narrative, Goldie calls
emplotment. Emplotment is a psychological process which can be more or
less successful. Emplotment can happen to those implicated in the
narrative, and also those who engage with the narrative. In other words,
narratives give those who are implicated by them and those who engage
with them, ways to consciously or subconsciously think about the events
mentally represented by that narrative.

The story your friend has told you can come with some emotional baggage.
This emotional baggage, in part, is what determines the appropriateness
of your response, whether you will take your usual way, or the other.
Assuming that being cited for a traffic violation equates to what most
of us think of as a negatively valenced experience, taking the longer
route would be the appropriate response in this situation.

\hypertarget{ii.-emotional-import}{%
\subsection{II. Emotional Import}\label{ii.-emotional-import}}

But when we think about this emotional import, it is not clear
initially. First, the proper response to a situation is outlined by the
emotional response to a given event. But often times, it is the
narrative that outlines the proper emotional response. Our friend may
say to us; ``oh man, you don't wanna go that way let me tell you . . .''
at which point he spins a yarn about narrowly being stopped and how if
it wasn't for the guy behind him . . . etc. As such

\begin{enumerate}
\def\labelenumi{\arabic{enumi}.}
\tightlist
\item
  Narratives first seem to outline the proper emotional response to a
  given event.
\end{enumerate}

But, they are also,

\begin{enumerate}
\def\labelenumi{\arabic{enumi}.}
\setcounter{enumi}{1}
\tightlist
\item
  distinct from emotions that are an agent's response to a given event
  or sequence of events.
\end{enumerate}

\begin{quote}
The emotional import of a narrative is thus distinct from any emotion
that might be felt in response: what happened was horrifying, and I, as
``external narrator,'' might or might not respond with horror to my
horrifying narrative of what happened. (\textbf{gold08?})
\end{quote}

The distinction then is between emotional import and emotional response.
We ought to contrast these two features.

\begin{enumerate}
\def\labelenumi{\arabic{enumi}.}
\item
  emotional import is an integral part of the narrative
\item
  emotional response is a response to it
\end{enumerate}

We may have felt the fear in our friend's voice. We saw him tremor as he
walked in the door. But while he tells the story, if he is even somewhat
adequate at doing so, it is us who feels the fear. We are transplanted
into the story where as the character, we become afraid.

As the internal character, it is I \emph{in the narrative} who feels the
emotion of fear. I respond the way someone who is afraid would in the
narrative, because in the narrative I find the incident to be
frightening. As an external narrator, I might also feel fear as a result
of conceiving this frightening incident, but the fear is not a necessary
response. In the movie \emph{Pulp Fiction}, while Vincent feels fear
having just accidentally killed Marvin, we as the external character
feel shock and confusion. Therefore this kind of response is external to
the narrative. Therefore, the agent as audience feels shock and as
character feel fear.

\hypertarget{second-example-where-they-differ}{%
\subparagraph{2 Second example, where they
differ}\label{second-example-where-they-differ}}

The agent as character may be humiliated by a more senior member at
work. In the moment, all they would feel is humiliation, embarrassment,
and shame. However, taking an external perspective on what happened, ``I
might conceive of what happened as deserving anger and resentment.''

By thinking through these narratives, we make the emotions appropriate
in response to the narrative. For instance, dramatic irony, the narrator
is not aware of crucial facts about what happened until later and as
such, the external emotion, the one in response to the internal emotion
is different from the internal one.

\hypertarget{third-example-variances}{%
\subparagraph{3 third example,
variances}\label{third-example-variances}}

\begin{itemize}
\item
  content of the thought-through narrative varies
\item
  as does one's grasp of what is happening in the actual world varies
\end{itemize}

Imagine a narrative scenario that explains current events. You might
have heard on the news that holiday spending is down from this time last
year. You can imagine that recent inflation may be the reason, or you
may imagine that a recently elected politician is. A similar example may
be that you are waiting for someone who has failed to arrive at the
agreed upon time. You might be waiting at the train station for someone
who should be there, but who has not yet shown up. You imagine her lying
injured somewhere, or deciding at the last minute to make other plans.

Each imagined scenario comes complete with its own set of emotions,
frustration, fear, compassion, etc.

\hypertarget{narrative-thinking-and-planning}{%
\subsection{Narrative Thinking and
Planning}\label{narrative-thinking-and-planning}}

When we imagine who we want to be, we may begin by thinking of what
steps we can actively take to become that person. However, sometimes,
once we set out on a path to becoming the person we desire, external
constraints can reveal to us the feasibility of our goals. Upon entering
university, we may have thought that we wanted a career in computer
science but realizing how difficult the math is, we may think it is not
worth it. Other factors can also constrain our choices.

Even better examples however, will concern our abilities to have close
meaningful relationships with others. For instance, we may desire to be
the sort of person whose ``door is always open.'' We might see ourselves
as someone who others seek out for advice, or companionship. However, we
can further imagine that our financial situations makes the thought of
inviting others over too uncomfortable. We may think that our house is
too small or our neighborhood too impoverished, rural or dangerous.

There are two kinds of factors, positive and negative ones. A negative
constraint is one wherein we pursue computer science but struggle with
the math and decide we would rather study philosophy. A positive
constraint however, is one wherein while we may be okay doing the math
required by the CS field, it would detract too much from what we really
liked about CS which is the analytic structure of argument. What we
liked in computer science however, was logic and not so much the math.
This is an example of a positive constraint. We discover some additional
desire that we have and this constrains our choice. The positive
constraint reveals an emerging goal or desire. This emerging goal or
desire presents new steps that we should take.

\begin{quote}
``Hypothetical imperatives are imperatives that tell you what you ought
to do if you want to achieve some stated goal; as Kant puts it, they
``declare a possible action to be practically necessary as a means to
the attainment of something else that one wills (or that one may will).
(\textbf{gold09?})
\end{quote}

The steps here being referenced are hypotheticals, they are imperatives
about what we should do given our stated goals. While hypotheticals are
implied by positive constraints, there is not initially a similar object
with respect to negative constraints. For instance, we may think that
negative constraints are associated not with \emph{oughts}, but with
\emph{ought-nots}, but it is not clear.

\begin{itemize}
\item
  Hypotheticals == imperatives about oughts in light of stated goals
\item
  When you fall short of these goals, you experience negative emotions.
\item
  the negative emotions motivate counterfactual thinking, thinking back
  to the crucial moment that contributed to your failure to realize your
  goals
\item
  Regret is an external emotion that arises in response to your decision
  underlying your goal failure
\item
  Each, negative emotions, counterfactual narrative thinking, and
  regrets can have a feedback effect on one's grasp of the related
  hypothetical imperatives, helping one to learn from one's mistakes.
\end{itemize}

\hypertarget{works-cited}{%
\section{Works Cited}\label{works-cited}}

\setlength{\parskip}{8pt}
\setlength{\parindent}{-0.15in}
\setlength{\leftskip}{0.15in}
\vspace*{-0.35in}

\noindent \singlespacing

\hypertarget{refs}{}
\begin{CSLReferences}{1}{0}
\leavevmode\vadjust pre{\hypertarget{ref-shus95}{}}%
Shusterman, Richard. 1995. {``Rap {Remix}: {Pragmatism},
{Postmodernism}, and {Other Issues} in the {House}.''} \emph{Critical
Inquiry} 22 (1): 150--58. \url{https://doi.org/10.1086/448785}.

\leavevmode\vadjust pre{\hypertarget{ref-shus07}{}}%
---------. 2007. {``Rap as {Art} and {Philosophy}.''} In, 419--28.
\url{https://doi.org/10.1002/9780470751640.ch29}.

\end{CSLReferences}

\end{document}
